\documentclass{article}

% CS230 project milestone template, built on the official NeurIPS 2026 style file.
% It uses the same setup as proposal_template/ and final_report_template/ so the
% proposal, milestone, and final report share one look. Keep the [final] option: it
% prints the author block and omits line numbers.
% natbib is loaded by the style file; this makes citations numeric, e.g. [1, 2].
% Remove the next line for author-year citations, e.g. (Vaswani et al., 2017).
\PassOptionsToPackage{numbers,compress}{natbib}
\usepackage[final]{neurips_2026}

% to avoid loading the natbib package, add option nonatbib:
%    \usepackage[final,nonatbib]{neurips_2026}

\usepackage[utf8]{inputenc} % allow utf-8 input
\usepackage[T1]{fontenc}    % use 8-bit T1 fonts
\usepackage{hyperref}       % hyperlinks
\usepackage{url}            % simple URL typesetting
\usepackage{booktabs}       % professional-quality tables
\usepackage{amsfonts}       % blackboard math symbols
\usepackage{nicefrac}       % compact symbols for 1/2, etc.
\usepackage{microtype}      % microtypography
\usepackage{xcolor}         % colors
\usepackage{graphicx}       % figures
% Page limit for the milestone, used in the text below. Update it here only.
\newcommand{\pagelimit}{3}

% Submission: upload this PDF on Gradescope.

% State the full title of your project. Keep the same title across the milestone and
% the final report where possible.
\title{Early Sepsis Prediction with Temporal and Graph Neural Networks}

% List the full name and SUNet ID of every team member.
% Use \And between authors to let LaTeX decide where to break the line, and
% \AND to force a line break at that point.
\author{%
  Nick Allen \\
  \texttt{nallen21@stanford.edu} \\
  \And
  Tyler Yamori-Little \\
  \texttt{tyamori@stanford.edu} \\
}

\begin{document}

\maketitle

\section{Introduction}
Sepsis requires rapid treatment, yet its early signs are dispersed across irregularly sampled ICU measurements. We study hourly early-warning predictions on the public PhysioNet/Computing in Cardiology Challenge 2019 data \cite{reyna2019}. At each hour, a model may use only measurements available by that hour and must decide whether to raise an alarm. The challenge's normalized utility rewards timely warnings, with the greatest benefit approximately six hours before onset, while penalizing false alarms and late detection. Our central question is whether trained temporal or physiological-variable representations improve this utility over a feature-engineered XGBoost baseline. We also plan to test whether reconstruction of masked observed values helps a graph model under severe measurement sparsity.

\section{Dataset}
The public labeled data contain 40,336 ICU stays from two hospitals (20,336 from A and 20,000 from B), comprising 1,552,210 hourly records. The team downloaded both labeled hospital cohorts and implemented patient-level splitting and causal input pipelines. The dataset is released under a Creative Commons Attribution 4.0 license. Each hour has 34 physiological variables and six demographic or contextual variables, plus \texttt{SepsisLabel}. Sepsis occurs in 7.27\% of stays and 1.80\% of hours; approximately 80.1\% of physiological entries are missing. Measurement practices also differ substantially by hospital.

We partitioned whole stays into train, development, and internal test sets of 28,235 / 6,053 / 6,048 patients (70/15/15), stratified by hospital, ever-septic status, and stay duration. No stay crosses partitions. The released label already begins six hours before the defined onset; we do not shift it again for ordinary training or scoring. For neural inputs, observed values are standardized with statistics fit on training patients. Missing values are forward-filled after a measurement and otherwise use a training-set median. Observation masks and elapsed time since each measurement remain explicit. The CNN and LSTM input pipelines use only current and past measurements.

\section{Approach}
Our reportable XGBoost baseline is a Sepsyd-style model built with XGBoost \cite{chen2016xgboost}: a five-hour stack of the 40 inputs and 34 physiological missingness indicators, with log transforms for selected right-skewed labs and forward filling. This is a retraining on our patient split, not the original challenge model or its hidden-test score.

We trained a three-layer causal convolutional network with 64 channels, kernel width 9, dropout 0.4, weight decay, and noise added to observed values during training. Its receptive field spans 25 hours. We also searched causal LSTMs and selected a three-layer, 64-hidden-unit model with a cosine learning-rate schedule. The neural models use Adam \cite{kingma2015adam} and class-weighted binary cross-entropy, $\mathcal{L}=-w_+y\log p-(1-y)\log(1-p)$; the training ratio gives a positive-hour weight of approximately 54.6. Model configurations and alarm thresholds were chosen using development utility. We additionally trained a causal attention model and an initial graph-attention model \cite{velickovic2018graph} over 34 physiological-variable nodes. The proposed masked-value reconstruction objective has not yet been implemented.

Normalized challenge utility is our primary metric; area under the precision--recall curve (AUPRC) is a secondary threshold-free metric. The alarm threshold is frozen after development-set selection. Our internal test set contains only hospitals A and B, so its scores are not directly comparable to published results on the challenge's different hidden cohort.

\paragraph{Remaining work.}
We will implement masked-value reconstruction: hide observed physiological values at training time, predict them from available graph features, and compare sepsis utility against the same architecture without reconstruction. The original observation mask will remain explicit so reconstructed values are not treated as measurements. We will also replace future-looking interpolation before using the exploratory tree model in a causal ensemble, assess errors and uncertainty by hospital, and refine graph temporal context. A strict Hospital A-to-B transfer experiment must fit preprocessing and graph statistics on A alone. The related CS224W work studies graph topology using shared data and infrastructure; this CS230 project focuses on trained temporal models and reconstruction.

\section{Partial Results}
Table~\ref{tab:results} summarizes selected runs on held-out patients. The Sepsyd-style XGBoost baseline currently outperforms the causal CNN and LSTM in utility, despite the latter establishing working deep-learning pipelines. The initial clinical graph model performs worse still, motivating stronger temporal features and the planned reconstruction objective. These negative results guide the next iteration.

\begin{table}[h]
\centering
\small
\begin{tabular}{lrrr}
\toprule
Model & Dev utility & Test utility & Test AUPRC\\
\midrule
Sepsyd-style XGBoost & 0.400 & 0.401 & 0.117\\
Regularized causal CNN & 0.374 & 0.385 & 0.099\\
Selected causal LSTM & 0.365 & 0.375 & 0.080\\
Selected causal attention & 0.354 & 0.371 & 0.086\\
Clinical graph attention & 0.308 & 0.300 & 0.072\\
Exploratory ensemble$^\dagger$ & 0.447 & 0.466 & 0.117\\
\bottomrule
\end{tabular}
\caption{Preliminary internal-holdout results, rounded from committed metrics. Multiple test results have now been inspected, so the test set is no longer a pristine final model-selection check.}
\label{tab:results}
\end{table}

\noindent $^\dagger$The ensemble is included for transparency, not as a valid causal early-warning result. Approximately 56\% of its weight is on a Separatrix-style tree model whose current linear interpolation can use later measurements to fill an earlier hour. The ensemble inherits this future-information leakage. Its apparent utility gain cannot establish real-time performance until that component is rebuilt with past-only features and reevaluated.

\section*{References}

% The list below is generated by BibTeX from references.bib; run
% "latexmk -pdf main.tex" (or pdflatex, bibtex, pdflatex, pdflatex).
\renewcommand{\bibsection}{}   % the "References" heading is already typeset above
\medskip
{\small
\bibliographystyle{plainnat}
\bibliography{references}
}

\end{document}
